% !TEX root = ../../../main.tex
% Block card: detection.impulse_events -- "Impulse events" (Part I, the blocks chapter).
\cardfield{What it does}
Turns a score that is mostly zero with occasional sharp spikes into events: one event per
spike, running from just before it to the next one.

\cardfield{When to use it}
Straight after \blockref{detection.deconvolve}, whose jump scores are exactly this kind of
score. It keeps a second drop that begins before the first has recovered as its own event,
which a plain cut on the signal would merge.

\cardfield{What goes in, what comes out}
In go the scores from the block before and the signal itself. Out comes one stretch of time
per event. Each event keeps where it starts (the last moment before the jump), where it goes
furthest from that level, and the total size of its jump.

\cardfield{Settings that matter}
\emph{threshold} first: a score above it is a jump. With Deconvolve before it, the score is in
noise widths, so 4 means ``a jump four times the noise''. You can drag it on the score curve.
\emph{merge\_within\_s} joins jumps closer than this, so a fall spread over a few samples is one
event. \emph{max\_extent\_s} caps how long an event may run when no next event ends it first.

\cardfield{Watch out for}
Neighbouring events overlap: each event starts \emph{pre\_s} seconds before its jump, so it
begins inside the previous one. An event's length is set by the next event, not by when the
trough recovers, so a lone event runs to the cap.
